\section{Reporting the Metric}

\subsection{The Curve Is the Metric}

The $a_k$ curve is the primary output of this framework.
It reveals the total learnable structure among the policy's outputs: the elbow location indicates when $\theta$ has learned most of the available patterns, $a_\infty$ reflects residual unpredictable variation, and $a_1$ reflects the unconditional surprise.
Overlaying curves from different policies on the same axes provides immediate visual diagnosis of how an intervention changes the output distribution, just as comparing scree plots does in PCA.

Everything below is a lossy summary of this curve.
We develop these summaries because curves are difficult to rank, average across prompt suites, or report in tables, but the reader should understand that each successive compression discards information.

\subsection{The \texorpdfstring{$(a_n, C, \sigma_\ell)$}{(a_n, C, sigma)} Triple}

The asymptotic floor $a_n$ (estimating $a_\infty$) and the per-byte coherence $C$ answer genuinely different questions:
\begin{itemize}[nosep]
    \item $a_n$ (bits/byte): ``How much surprise per byte remains after $\theta$ has seen the other responses?''
    \item $C$ (dimensionless, in $(0,1)$): ``How likely, per byte, does $\theta$ find $\pi$'s outputs?''
\end{itemize}

This separation enables diagnostic reasoning.
Comparing policy $\pi_1$ against $\pi_2$, one might observe that $\pi_2$ has lower $a_n$ (its responses become easier to predict once $\theta$ has seen others) while its $C$ remains similar to $\pi_1$'s (each output is individually plausible).
Alternatively, $\pi_2$ might have lower $C$ (its outputs are less plausible to $\theta$) without a corresponding change in $a_n$.
These are qualitatively different failure modes that a single scalar would conflate.

Because $C$ is a mean, it can mask heterogeneity.
We therefore pair it with the coherence spread $\sigma_\ell$ (Section~\ref{sec:coherence-hetero}), giving the standard summary triple $(a_n, C, \sigma_\ell)$.
Note that $\sigma_\ell$ is defined in per-byte units (the standard deviation of per-byte cross-entropies $h_\theta(r_i \mid p)$) rather than total bits, because this isolates quality variation from length variation, making it a better diagnostic of coherence heterogeneity.

\subsection{Scalar Summary: Plausibility-Weighted Diversity}\label{sec:scalar}

When a single number is required (e.g., for ranking policies in a table), we combine coherence with the asymptotic surprise floor.
Diverse responses remain surprising even after full conditioning ($a_\infty$ high), while paraphrases become predictable ($a_\infty$ low).
This motivates:
\begin{equation}\label{eq:D-ainf}
    D_{Ca_n} = C \times a_n
\end{equation}
measuring ``how many bits of surprise remain per byte after $\theta$ has learned what it can from the other responses, weighted by output plausibility.''

In practice, $a_\infty$ is estimated by the last observed point $a_n$ on the permutation-averaged per-byte curve.
Section~\ref{sec:tevet} compares $D_{Ca_n}$ against alternative scalars (raw $a_n$, $C$ alone, the excess-entropy-based $C \times E$) and finds that $D_{Ca_n}$ outperforms all of them on external benchmarks.

The intended edge cases on the five validation scenarios (Section~\ref{sec:scenario-validation}) behave correctly:
\begin{enumerate}[nosep]
    \item \textbf{Pure noise}: $C$ low (incoherent), $a_\infty$ high, so $D_{Ca_n}$ is low (suppressed by $C$).
    \item \textbf{Multi-mode incoherent}: $C$ low, $a_\infty$ high, so $D_{Ca_n}$ is low (suppressed by $C$).
    \item \textbf{Multi-mode coherent}: $C$ high, $a_\infty$ high, so $D_{Ca_n}$ is high (the intended winner).
    \item \textbf{One-mode}: $C$ high, $a_\infty$ low ($\theta$ learns the mode), so $D_{Ca_n}$ is low (suppressed by $a_\infty$).
    \item \textbf{Mixed}: $C$ mid, $a_\infty$ high, so $D_{Ca_n}$ is mid.
\end{enumerate}

\paragraph{Typical range of $C$.} In our Tevet evaluation with Qwen2.5-3B, the central 90\% (5th to 95th percentile) of per-response bits-per-byte values across \tevetCoherenceRespCount{} responses span $[\tevetCoherenceBpbLow, \tevetCoherenceBpbHigh]$, giving $C \in [\tevetCoherenceCLow, \tevetCoherenceCHigh]$; the mean per-set $C$ across \tevetCoherenceSetCount{} response sets is \tevetCoherenceMeanC.
If coherence variation dominates diversity variation in a particular application, $C$ can be normalized (e.g., dividing by the coherence of $\theta$'s own samples); see Appendix~\ref{app:c-normalization} for details.

\subsection{Coherence Heterogeneity}\label{sec:coherence-hetero}

The coherence $C$ summarizes the plausibility of $\pi$'s outputs as a single number, but a policy may mix coherent and incoherent modes.
A policy with 5 coherent modes and 3 incoherent modes will have a high $a_n$ (since the modes are genuinely distinct and $\theta$ cannot fully predict any one from the others) while $C$ is dragged down by the incoherent samples.
The pair $(a_n, C)$ cannot distinguish ``8 modes, 3 of which are junk'' from ``8 uniformly mediocre modes.''

To diagnose this, we report the \textbf{coherence spread} $\sigma_\ell$, defined as the standard deviation of the per-response cross-entropies:
\begin{equation}
    \sigma_\ell = \mathrm{std}\bigl(\{h_\theta(r_i \mid p)\}_{i=1}^n\bigr)
\end{equation}

Since $C = 2^{-\bar{\ell}}$ where $\bar{\ell} = \frac{1}{n}\sum_i h_\theta(r_i \mid p)$, the spread $\sigma_\ell$ induces a multiplicative coherence interval.
Define:
\begin{equation}
    C_+ = 2^{-(\bar{\ell} - \sigma_\ell)} = C \cdot 2^{\sigma_\ell}, \qquad C_- = 2^{-(\bar{\ell} + \sigma_\ell)} = C \cdot 2^{-\sigma_\ell}
\end{equation}

$C_+$ is the coherence of the better-than-average responses and $C_-$ is the coherence of the worse-than-average responses.
These yield a \textbf{diversity uncertainty band}:
\begin{equation}
    D_\pm = C_\pm \times a_n
\end{equation}

The width of this band, $D_+ - D_- = C(2^{\sigma_\ell} - 2^{-\sigma_\ell}) \times a_n$, is zero when all responses have equal coherence and large when the policy mixes coherent and incoherent modes.
A wide band is the diagnostic signal for heterogeneous output quality.

\subsection{What to Report}

In order of informativeness:
\begin{enumerate}[nosep]
    \item The $a_k$ curve, the primary output.
It reveals the total learnable structure, its rate of acquisition, residual variation ($a_\infty$), and unconditional surprise ($a_1$).
When comparing policies, overlaying curves on the same axes provides immediate visual insight.
    \item The summary triple $(a_n, C, \sigma_\ell)$ together with the scalar $D_{Ca_n} = C \times a_n$ (bits/byte).
When $\sigma_\ell$ is small, $D_{Ca_n}$ is a reliable scalar; when $\sigma_\ell$ is large, report the diversity band $[D_-, D_+]$ and examine which modes are incoherent via the per-response $c_i = 2^{-h_\theta(r_i \mid p)}$ values.
\end{enumerate}

The score is defined per prompt $p$.
When comparing policies across a prompt suite, one can average $D_{Ca_n}$ over prompts or take differences; see Appendix~\ref{app:aggregation} for details.

